% !PRINTED 309
\weberpage{309}
The following remark is still of interest here. The resolvents of degree $p$, whose roots are the quantities (9), (10), (11), still contain $\sqrt{\pm p}$ in their coefficients, according to §\,78.

But the group $\mathfrak{L}_0$ is extended to the group $\mathfrak{L}$ if we add \gesperrt{a} linear substitution of the form §\,78, (13) in which $ad - bc$ is a quadratic non-residue, thus for instance: % Source scan clearly prints ad - bc.
\[
\begin{aligned}
\text{for } p&=5 & &(z,\,2z)\\
\text{for } p&=7,\,11 & &(z,\,-z).
\end{aligned}
\]
By applying this substitution, however, the value (9)
\[
V = (v_\infty - v_0)(v_1 - v_4)(v_2 - v_3)
\]
passes into its opposite, and from this it follows that, when $p=5$, $V^2$ satisfies an equation of degree 5 with \gesperrt{rational} numerical coefficients. From this one concludes that in the equation for $V$ the coefficients of the odd powers have the factor $\sqrt5$, while the others are rational; or that for the unknown $\sqrt5\,V$ one obtains an equation with \gesperrt{rational coefficients}.

For $p=7,\,11$ the expressions corresponding to the two signs in (10) or (11) pass into one another through the interchange $(z,-z)$. The resolvents of degree 7 or 11, which one of the two expressions (10) or (11) respectively satisfies, therefore pass into one another through a change of the sign of $\sqrt{-7}$, $\sqrt{-11}$.

§\,83. \gesperrt{Various Resolvents of Degree 5 for the 5th Degree of Transformation.}

In the formation of the resolvents of degree $p$ we restrict ourselves here to the case $p=5$.

These resolvents can be given very manifold forms; for one may not only choose, in the function $V$, (9) of the preceding section, the roots of an \gesperrt{arbitrary} transformation equation for the quantities $v$, but may also replace $V$ by various other functions, for example by
\[
(v_\infty + v_0)(v_1+v_4)(v_2+v_3),
\]
or by
\[
v_\infty v_0 + v_1 v_4 + v_2 v_3.
\]
